\chapter{برنامه‌های کوچک کاربردی}
\label{ch:projects}

\begin{goals}
\begin{itemize}
  \item کنار هم گذاشتن همه‌ی آموخته‌ها در چند برنامه‌ی کامل
  \item شکستن یک مسئله‌ی روزمره به قدم‌های کوچک
  \item خواندن یک برنامه‌ی چند تابعی و فهمیدن جریان آن
  \item ایده‌هایی برای ادامه‌ی راه
\end{itemize}
\end{goals}

\section{از ابزار تا برنامه}

در فصل‌های گذشته ابزارهای پایه‌ی برنامه‌نویسی را یکی‌یکی شناختید: چاپ،
متغیر، نوع‌ها، محاسبه، شرط، حلقه، تابع، آرایه و رشته. هر برنامه‌ای که در
جهان نوشته شده، هر قدر بزرگ، از همین ابزارها ساخته شده است. در این فصل
چند برنامه‌ی کوچک اما کامل می‌نویسیم؛ برنامه‌هایی که شبیه کارهای واقعی‌اند
و در هر کدام چند ابزار با هم کار می‌کنند.

برای هر برنامه اول مسئله را می‌گوییم، بعد آن را به قدم‌های کوچک می‌شکنیم و
بعد برنامه را می‌نویسیم. پیشنهاد می‌کنیم پیش از خواندن برنامه‌ی ما، خودتان
تلاش کنید آن را بنویسید؛ سپس دو نسخه را مقایسه کنید. هیچ اشکالی ندارد که
راه شما با راه ما فرق داشته باشد.

\section{قبض برق}

\textbf{مسئله.} شرکت برق مصرف ماهانه را پلکانی حساب می‌کند: ۱۰۰ کیلووات‌ساعت
اول هر واحد ۵۰۰ تومان، ۲۰۰ واحد بعدی هر واحد ۱٬۰۰۰ تومان و بیش از آن هر واحد
۲٬۰۰۰ تومان. مبلغ قبض چند مشترک را حساب کنید.

\textbf{قدم‌ها.} تابعی می‌خواهیم که مصرف را بگیرد و مبلغ را برگرداند. مصرف
را پله‌پله کم می‌کنیم: اول تا ۱۰۰ واحد را با نرخ اول، بعد تا ۲۰۰ واحد را با
نرخ دوم و بقیه را با نرخ سوم. سپس برای آرایه‌ای از مصرف‌ها، مبلغ هر یک را
چاپ می‌کنیم.

\program{قبض برق پلکانی}
\begin{salam}
تابع مبلغ قبض(مصرف: صحیح): صحیح:
    متغیر باقی := مصرف
    متغیر مبلغ := 0

    // پله‌ی اول: تا ۱۰۰ واحد، هر واحد ۵۰۰ تومان
    پله یک := باقی > 100 ؟ 100 : باقی
    مبلغ += پله یک * 500
    باقی -= پله یک

    // پله‌ی دوم: تا ۲۰۰ واحد بعدی، هر واحد ۱۰۰۰ تومان
    پله دو := باقی > 200 ؟ 200 : باقی
    مبلغ += پله دو * 1000
    باقی -= پله دو

    // پله‌ی سوم: هر چه ماند، هر واحد ۲۰۰۰ تومان
    مبلغ += باقی * 2000
    بازگشت مبلغ
تمام

تابع اصلی:
    مصرف مشترکان := [80, 100, 250, 400]
    هر مصرف در مصرف مشترکان:
        چاپ "مصرف", مصرف, "واحد:", مبلغ قبض(مصرف), "تومان"
    تمام
تمام
\end{salam}

\begin{outputfa}
مصرف 80 واحد: 40000 تومان
مصرف 100 واحد: 50000 تومان
مصرف 250 واحد: 200000 تومان
مصرف 400 واحد: 450000 تومان
\end{outputfa}

مصرف ۲۵۰ واحد را دنبال کنید: ۱۰۰ واحد در پله‌ی اول (۵۰ هزار)، ۱۵۰ واحد در
پله‌ی دوم (۱۵۰ هزار)، و چیزی برای پله‌ی سوم نمی‌ماند؛ جمعاً ۲۰۰ هزار تومان.
عملگر \code{?:} در هر پله می‌گوید «همه‌ی ظرفیت پله، یا هر چه مانده، هر کدام
که کمتر است».

\begin{tip}[ایده برای ادامه]
نرخ‌ها را به \fcode{ثابت} تبدیل کنید و پله‌ی چهارمی اضافه کنید. سپس در
تابع اصلی جمع کل درآمد شرکت را هم حساب کنید.
\end{tip}

\section{نمودار نمره‌ها}

\textbf{مسئله.} نمره‌های یک کلاس را داریم و می‌خواهیم ببینیم چند نفر در هر
بازه (زیر ۱۰، ۱۰ تا ۱۴، ۱۵ تا ۱۷، ۱۸ تا ۲۰) هستند و این را به شکل یک
نمودار ستاره‌ای نشان دهیم.

\textbf{قدم‌ها.} چهار شمارنده لازم داریم؛ برای هر نمره تصمیم می‌گیریم که به
کدام شمارنده اضافه شود. سپس برای هر بازه، یک خط چاپ می‌کنیم که به تعداد
شمارنده ستاره دارد. چاپ ستاره‌ها را به یک تابع می‌سپاریم.

\program{نمودار ستاره‌ای نمره‌ها}
\begin{salam}
تابع سطر نمودار(عنوان: رشته, تعداد: صحیح):
    متغیر ستاره ها := ""
    تکرار تعداد:
        ستاره ها += "*"
    تمام
    چاپ عنوان + ":", ستاره ها, "(" + تعداد + ")"
تمام

تابع اصلی:
    نمرات := [12, 19, 8, 15, 20, 17, 11, 9, 16, 18, 14, 13]
    متغیر مردود := 0
    متغیر متوسط := 0
    متغیر خوب := 0
    متغیر عالی := 0

    هر نمره در نمرات:
        اگر نمره < 10: مردود++
        وگرنه نمره < 15: متوسط++
        وگرنه نمره < 18: خوب++
        وگرنه: عالی++
        تمام
    تمام

    چاپ "نمودار نمره‌های کلاس"
    سطر نمودار("زیر ۱۰   ", مردود)
    سطر نمودار("۱۰ تا ۱۴ ", متوسط)
    سطر نمودار("۱۵ تا ۱۷ ", خوب)
    سطر نمودار("۱۸ تا ۲۰ ", عالی)
تمام
\end{salam}

\begin{outputfa}
نمودار نمره‌های کلاس
زیر ۱۰   : ** (2)
۱۰ تا ۱۴ : **** (4)
۱۵ تا ۱۷ : *** (3)
۱۸ تا ۲۰ : *** (3)
\end{outputfa}

به شکل فشرده‌ی زنجیره‌ی شرط دقت کنید: هر شاخه یک دستور کوتاه دارد و در
همان خط نوشته شده و فقط یک \fcode{تمام} کل زنجیره را می‌بندد. برای
تراز شدن ستون ستاره‌ها، به عنوان‌ها فاصله‌ی اضافی چسبانده‌ایم.

\begin{tip}[ایده برای ادامه]
به جای چهار شمارنده‌ی جداگانه، یک آرایه‌ی \fcode{متغیر} با چهار صفر بسازید
و با اندیس، شمارنده‌ی مناسب را افزایش دهید. سپس با یک حلقه‌ی \fcode{هر} و
یک آرایه‌ی هم‌راستا از عنوان‌ها، نمودار را چاپ کنید.
\end{tip}

\section{صورت‌حساب فروشگاه}

\textbf{مسئله.} فهرست کالاهای سبد خرید یک مشتری را داریم: نام، قیمت واحد و
تعداد. باید صورت‌حسابی چاپ کنیم که مبلغ هر ردیف، جمع کل، تخفیف (اگر خرید
از ۵۰۰ هزار تومان بیشتر بود ۱۰ درصد) و مبلغ نهایی را نشان دهد.

\textbf{قدم‌ها.} سه آرایه‌ی هم‌راستا داریم. با یک حلقه روی نام‌ها می‌رویم،
مبلغ هر ردیف را حساب و چاپ می‌کنیم و به جمع کل می‌افزاییم. بعد از حلقه،
تخفیف را با یک شرط تعیین می‌کنیم.

\program{صورت‌حساب فروشگاه}
\begin{salam}
ثابت آستانه := 500000
ثابت تخفیف := 10

تابع اصلی:
    کالاها := ["برنج", "روغن", "چای", "شکر"]
    قیمتها := [180000, 95000, 120000, 40000]
    تعدادها := [2, 1, 1, 3]
    متغیر جمع := 0

    چاپ "صورت‌حساب"
    چاپ "-----------------------------"
    هر (شماره, کالا) در کالاها:
        مبلغ := قیمتها[شماره] * تعدادها[شماره]
        چاپ کالا, "x", تعدادها[شماره], "=", مبلغ
        جمع += مبلغ
    تمام
    چاپ "-----------------------------"
    چاپ "جمع کل:", جمع

    اگر جمع > آستانه:
        مبلغ تخفیف := جمع * تخفیف / 100
        چاپ "تخفیف", تخفیف, "درصد:", مبلغ تخفیف
        چاپ "قابل پرداخت:", جمع - مبلغ تخفیف
    وگرنه:
        چاپ "قابل پرداخت:", جمع
        چاپ "با", آستانه - جمع, "تومان خرید بیشتر، ۱۰ درصد تخفیف بگیرید"
    تمام
تمام
\end{salam}

\begin{outputfa}
صورت‌حساب
-----------------------------
برنج x 2 = 360000
روغن x 1 = 95000
چای x 1 = 120000
شکر x 3 = 120000
-----------------------------
جمع کل: 695000
تخفیف 10 درصد: 69500
قابل پرداخت: 625500
\end{outputfa}

تعداد شکر را به ۱ تغییر دهید تا شاخه‌ی \fcode{وگرنه} را هم ببینید. متغیر
\fcode{مبلغ} درون حلقه و بدون \fcode{متغیر} تعریف شده، چون در هر دور یک
بار مقدار می‌گیرد و همان است؛ فقط \fcode{جمع} است که در طول حلقه تغییر
می‌کند.

\begin{tip}[ایده برای ادامه]
یک آرایه‌ی چهارم برای «کالای تخفیف‌دار» (منطقی) اضافه کنید که روی آن
کالاها ۵ درصد تخفیف جداگانه اعمال شود. یا گران‌ترین ردیف صورت‌حساب را
پیدا کنید و در پایان چاپ کنید.
\end{tip}

\section{تقویم یک ماه}

\textbf{مسئله.} می‌خواهیم تقویم یک ماه را چاپ کنیم: هفت ستون از شنبه تا
جمعه، با روزهای ماه در جای درست. برای این کار باید بدانیم ماه چند روز
دارد و روز اول آن چه روزی از هفته است (۰ برای شنبه تا ۶ برای جمعه).

\textbf{قدم‌ها.} اول سرستون‌ها را چاپ می‌کنیم. بعد یک رشته می‌سازیم و به
تعداد روزهایی که پیش از روز اول ماه در هفته گذشته، جای خالی می‌گذاریم.
سپس روزها را یکی‌یکی می‌چسبانیم و هر بار که به جمعه رسیدیم، خط را چاپ
می‌کنیم و خط تازه‌ای شروع می‌کنیم. برای این که ستون‌ها تراز بمانند،
عددهای یک‌رقمی را با یک فاصله‌ی اضافی چاپ می‌کنیم.

\program{تقویم ماه}
\begin{salam}
تابع دو رقمی(عدد: صحیح): رشته:
    اگر عدد < 10: بازگشت " " + عدد تمام
    بازگشت "" + عدد
تمام

تابع اصلی:
    تعداد روز := 31
    روز اول := 3     // ۰ شنبه ... ۶ جمعه؛ این ماه از سه‌شنبه شروع می‌شود

    چاپ "شنبه یکشنبه دوشنبه سه‌شنبه چهارشنبه پنجشنبه جمعه"
    چاپ " ش   ی   د   س   چ   پ   ج"
    متغیر خط := ""
    متغیر ستون := 0

    // جاهای خالی پیش از روز اول
    تکرار روز اول:
        خط += "    "
        ستون++
    تمام

    تکرار 1 تا تعداد روز با روز:
        خط += دو رقمی(روز) + "  "
        ستون++
        اگر ستون == 7:
            چاپ خط
            خط = ""
            ستون = 0
        تمام
    تمام
    اگر ستون > 0: چاپ خط تمام
تمام
\end{salam}

\begin{outputfa}
شنبه یکشنبه دوشنبه سه‌شنبه چهارشنبه پنجشنبه جمعه
 ش   ی   د   س   چ   پ   ج
             1   2   3   4
 5   6   7   8   9  10  11
12  13  14  15  16  17  18
19  20  21  22  23  24  25
26  27  28  29  30  31
\end{outputfa}

شرط پایانی \fcode{اگر ستون \lr{>} 0} برای این است که هفته‌ی ناتمام آخر هم
چاپ شود؛ بدون آن، روزهای ۲۶ تا ۳۱ در \fcode{خط} می‌ماندند و هرگز چاپ
نمی‌شدند. این نوع «کار ناتمام بعد از حلقه» یکی از جاهایی است که
برنامه‌نویسان زیاد فراموش می‌کنند.

\begin{tip}[ایده برای ادامه]
با \fcode{تطبیق}، از روی شماره‌ی ماه تعداد روزهایش را به دست آورید، و روز
اول هر ماه را از روز اول ماه پیش حساب کنید (باقی‌مانده‌ی جمع بر ۷). آن وقت
می‌توانید تقویم کل سال را چاپ کنید.
\end{tip}

\section{تحلیل یک جمله}

\textbf{مسئله.} جمله‌ای داریم و می‌خواهیم بدانیم چند واژه دارد، درازترین
واژه‌اش کدام است و یک واژه‌ی مشخص چند بار در آن آمده است.

\textbf{قدم‌ها.} جمله را با فاصله می‌شکافیم. با یک حلقه روی واژه‌ها
می‌رویم: تعداد را می‌شماریم، درازترین را با الگوی «بیشترین» پیدا می‌کنیم و
هر جا واژه با واژه‌ی مورد نظر برابر بود، شمارنده‌ای را افزایش می‌دهیم.

\program{تحلیل جمله}
\begin{salam}
تابع اصلی:
    جمله := "سلام زبان برنامه‌نویسی است و سلام یعنی درود"
    واژه هدف := "سلام"

    واژه ها := جمله.بشکاف(" ")
    متغیر درازترین := ""
    متغیر شمار هدف := 0

    تکرار واژه ها.طول() با شماره:
        واژه := واژه ها.بگیر(شماره)
        اگر واژه.طول() > درازترین.طول():
            درازترین = واژه
        تمام
        اگر واژه == واژه هدف:
            شمار هدف++
        تمام
    تمام

    چاپ "تعداد واژه‌ها:", واژه ها.طول()
    چاپ "درازترین واژه:", درازترین
    چاپ "واژه‌ی «" + واژه هدف + "»", شمار هدف, "بار آمده است"
تمام
\end{salam}

\begin{outputfa}
تعداد واژه‌ها: 8
درازترین واژه: برنامه‌نویسی
واژه‌ی «سلام» 2 بار آمده است
\end{outputfa}

\fcode{درازترین} با رشته‌ی خالی شروع می‌شود که طولش صفر است؛ پس اولین
واژه حتماً جایش را می‌گیرد و از آن به بعد فقط واژه‌های درازتر. همان الگوی
«بیشترین» است، این بار برای رشته‌ها.

\begin{tip}[ایده برای ادامه]
کوتاه‌ترین واژه را هم پیدا کنید. سپس جمله‌ی تازه‌ای بسازید که در آن هر
واژه‌ی هدف با واژه‌ی دیگری (مثلاً «درود») جایگزین شده باشد.
\end{tip}

\section{و بعد از این؟}

اگر تا این‌جا آمده‌اید و برنامه‌های این فصل را فهمیده‌اید، دیگر
برنامه‌نویس هستید؛ شاید تازه‌کار، اما برنامه‌نویس. از این‌جا به بعد چند
راه پیش رو دارید:

\begin{itemize}
  \item \textbf{بیشتر بنویسید.} هر کار کوچک روزمره را که با عدد یا متن
        سروکار دارد، به برنامه تبدیل کنید: محاسبه‌ی سود وام، برنامه‌ی
        هفتگی درس، آمار مصرف آب خانه. تمرین‌های پایان فصل‌ها را دوباره
        ببینید و آن‌هایی را که رد کرده‌اید بنویسید.
  \item \textbf{برنامه‌های نمونه‌ی ویرایشگر برخط را بخوانید.} حالا همه‌ی
        آن‌ها را می‌فهمید. بعضی از آن‌ها چیزهایی را نشان می‌دهند که در این
        کتاب نبود، مثل \fcode{ساختار} برای دسته‌بندی داده‌های مرتبط، یا
        \fcode{وکتور} برای فهرست‌هایی که بزرگ و کوچک می‌شوند.
  \item \textbf{کتاب‌های بعدی سلام.} این کتاب پایه بود. سلام امکانات
        بسیار بیشتری دارد: کتابخانه‌ای بزرگ از ابزارهای آماده، ساختن
        صفحه‌های وب، کار با پرونده‌ها و شبکه، و بسیاری چیزهای دیگر. همه‌ی
        این‌ها بر همان پایه‌ای بنا می‌شوند که این‌جا آموختید.
\end{itemize}

مهم‌تر از همه: از اشتباه نترسید. هر پیام خطایی که می‌بینید، یک درس کوچک
است؛ و هر برنامه‌ای که بالاخره درست کار می‌کند، یک پیروزی کوچک. لذت
برنامه‌نویسی در همین پیروزی‌های کوچک است.

\begin{summary}
\begin{itemize}
  \item هر مسئله را به قدم‌های کوچک بشکنید و برای هر قدم روشن، یک تابع
        یا یک بلوک بنویسید.
  \item الگوهایی که آموختید (شمارنده، جمع، بیشترین، جست‌وجو، ساختن رشته
        در حلقه) در برنامه‌های واقعی بارها تکرار می‌شوند.
  \item بعد از هر حلقه بپرسید: آیا کار ناتمامی مانده که باید انجام شود؟
  \item برنامه‌ای که کار می‌کند نقطه‌ی پایان نیست؛ آن را بهتر، کوتاه‌تر و
        خواناتر کنید.
\end{itemize}
\end{summary}

\section*{تمرین‌ها}

\exercise برنامه‌ای بنویسید که برای یک وام با مبلغ، نرخ سود سالانه و مدت
(به ماه)، قسط ماهانه‌ی ساده (مبلغ به علاوه‌ی سود کل، تقسیم بر تعداد ماه)
را حساب کند و جدول ماه‌به‌ماه مانده‌ی بدهی را چاپ کند.

\exercise دمای هر روز یک ماه را در آرایه‌ای بگذارید و گزارشی چاپ کنید
شامل: میانگین ماه، گرم‌ترین و سردترین روز، تعداد روزهای بالای ۳۰ درجه و
طولانی‌ترین رشته‌ی روزهای پیاپی بالای ۳۰ درجه.

\exercise فهرستی از نام‌ها و ساعت‌های کار هفتگی کارکنان را در دو آرایه‌ی
هم‌راستا بگذارید. با نرخ ساعتی ثابت، دستمزد هر نفر را حساب کنید؛ ساعت‌های
بیش از ۴۰ در هفته، اضافه‌کاری با نرخ یک‌ونیم برابر حساب می‌شود. جمع کل
دستمزدها را هم چاپ کنید.

\exercise یک «تبدیل‌گر واحد» بنویسید: تابعی که مقداری و نام واحد مبدأ و
مقصد (مثلاً \code{"km"} و \code{"m"}) را بگیرد و با \fcode{تطبیق} تبدیل
را انجام دهد. دست‌کم پنج جفت واحد را پشتیبانی کنید.

\exercise یک متن چندخطی (با \code{\textbackslash n}) را در رشته‌ای بگذارید،
آن را با \fcode{بشکاف("\textbackslash n")} به خط‌ها بشکافید و برای هر خط،
شماره‌ی خط و تعداد واژه‌هایش را چاپ کنید. در پایان بگویید کدام خط بیشترین
واژه را دارد.
